% Part: first-order-logic
% Chapter: introduction
% Section: models-theories

\documentclass[../../../include/open-logic-section]{subfiles}

\begin{document}

\olfileid{fol}{int}{mod}

\olsection{النماذج والنظريات}

متى عرّفنا نحو منطق الرتبة الأولى ودلالته، أمكننا أن نشرع في بحث خصائص
ال!!{structure}s والمفاهيم الدلالية. ويمكننا أيضًا أن نعرّف أنساق
ال!!{derivation} ونبحثها. ويمكننا أن نسأل عن مجموعة من
ال!!{sentence}s: أي ال!!{structure}s تجعل جميع ال!!{sentence}s في
تلك المجموعة صادقة؟ إذا أعطيت مجموعة !!{sentence}s~$\Gamma$، سميت
ال!!^a{structure}~$\Struct{M}$ التي ترضيها \emph{نموذجًا
لـ~$\Gamma$}. وقد نبدأ من~$\Gamma$ ونحاول إيجاد نماذجها: فما صورتها؟
وكم ينبغي أن تكون كبيرة أو صغيرة؟ لكن قد نبدأ أيضًا بـ!!{structure}
واحدة أو بمجموعة من ال!!{structure}s ونسأل: أي ال!!{sentence}s تصدق
فيها؟ وهل توجد !!{sentence}s \emph{تميز} هذه ال!!{structure}s، بمعنى
أنها تصدق في تلك البنى دون غيرها؟ وهذه الأنواع من الأسئلة هي ميدان
\emph{نظرية النماذج}. وهي أيضًا أساس \emph{المنهج البديهي}: وصف
مجموعة من ال!!{structure}s بمجموعة من ال!!{sentence}s، هي بديهيات
نظرية ما. ويغدو ذلك ممكنًا بفضل ملاحظة أن ال!!{sentence}s التي تلزم في منطق
الرتبة الأولى عن البديهيات هي بالضبط الصادقة في جميع نماذجها.

ولنتأمل، مثالًا بالغ البساطة، الترتيبات المسبقة. فالترتيب المسبق
علاقة~$R$ على مجموعة ما~$A$ تكون انعكاسية ومتعدية معًا. والمجموعة~$A$
المزودة بالعلاقة الثنائية $R \subseteq A \times A$ هي بالضبط ما
نحتاج إليه لإعطاء !!a{structure} للغة من لغات الرتبة الأولى ذات رمز
علاقة ثنائي واحد~$\Obj P$: إذ نجعل
$\Domain{M} = A$ و$\Assign{\Obj P}{M} = R$. ولما كانت $R$ ترتيبًا
مسبقًا، فهي انعكاسية ومتعدية، ويمكننا إيجاد مجموعة~$\Gamma$ من
ال!!{sentence}s في منطق الرتبة الأولى تقول ذلك:
\begin{align*}
  & \lforall[\Obj v_0][\Atom{\Obj P}{\Obj v_0,\Obj v_0}]\\
  & \lforall[\Obj v_0][\lforall[\Obj v_1][\lforall[\Obj v_2][((\Atom{\Obj P}{\Obj v_0,\Obj v_1} \land \Atom{\Obj P}{\Obj v_1,\Obj v_2}) \lif \Atom{\Obj P}{\Obj v_0,\Obj v_2})]]]
\end{align*}
وهاتان ال!!{sentence}s ليستا إلا ترميزًا لعبارتي «لكل~$x$،
$Rxx$» ($R$ انعكاسية) و«إذا كان $Rxy$ و$Ryz$ فإن $Rxz$ أيضًا»
($R$ متعدية). ونرى أن ال!!^a{structure}~$\Struct{M}$ نموذج لهاتين
ال!!{sentence}s~$\Gamma$ إذا وفقط إذا كانت $R$ (أي
$\Assign{\Obj P}{M}$) ترتيبًا مسبقًا على~$A$ (أي على
$\Domain{M}$). وبعبارة أخرى، نماذج $\Gamma$ هي الترتيبات المسبقة
بالضبط. وكل خاصية لجميع الترتيبات المسبقة يمكن التعبير عنها في لغة
الرتبة الأولى التي لا !!{predicate} فيها غير~$\Obj P$ (كالانعكاسية
والتعدي أعلاه) تلزم عن ال!!{sentence}s الاثنتين في~$\Gamma$، والعكس
صحيح. ولذلك فكل ما نبرهنه عن نماذج~$\Gamma$ نكون قد برهناه عن جميع
الترتيبات المسبقة.

وستكون لكل نظرية معينة وفئة من النماذج (مثل $\Gamma$ وجميع الترتيبات
المسبقة) أسئلة مهمة عما يمكن التعبير عنه في لغة الرتبة الأولى المقابلة
وما لا يمكن التعبير عنه فيها. وثمة خصائص لل!!{structure}s مهمة في
جميع اللغات وفئات النماذج، وهي التي تتعلق بحجم ال!!{domain}. فيمكن
دائمًا التعبير، مثلًا، عن أن ال!!{domain} يحتوي على $n$ من
ال!!{element}s بالضبط، لأي~$n \in \PosInt$. ويمكن أيضًا، باستعمال
مجموعة لا نهائية من ال!!{sentence}s، التعبير عن أن ال!!{domain} لا
نهائي. لكن لا يمكن التعبير عن أن المجال متناه، أو عن أن مجموعة
المجال !!{nonenumerable}. وهذه النتائج المتعلقة بقيود لغات الرتبة الأولى من
نتائج مبرهنتي التراص ولوفنهايم--سكولم.

\end{document}
